% Бүлэг 2: Системийн шинжилгээ

\chapter{Системийн шинжилгээ ба шаардлагын тодорхойлолт}
\label{Chapter2}

\section{Системийн ерөнхий зорилго ба үйл ажиллагаа}
Энэхүү систем нь дижитал хөгжим үйлдвэрлэлийн явцад эхлэн суралцагчдад тулгардаг техникийн саад бэрхшээлийг хиймэл оюун ухааны тусламжтайгаар шийдвэрлэх зорилготой. Систем нь хэрэглэгчийн бүтээсэн аудио файлд спектраль анализ хийж, миксинг (mixing) зөвлөмж өгөх, MusicVAE загвараар шинэ аяны санаа үүсгэх, болон RAG технологиор баталгаажсан онолын мэдлэг олгох чатботоос бүрдэнэ.

\section{Системийн оролцогчид (Actors)}
Системийг ашиглах үндсэн оролцогчдыг дараах байдлаар тодорхойлов:
\begin{itemize}
    \item \textbf{Хэрэглэгч (Продюсер):} Өөрийн аудио болон MIDI файлыг системд байршуулж, AI анализ болон зөвлөмж авах үндсэн оролцогч.
    \item \textbf{Мэргэжилтэн (Mentor):} Системд сургалтын материал оруулах, баталгаажсан өгөгдлийн санг баяжуулах продюсерууд.
    \item \textbf{Админ:} Системийн хэвийн ажиллагаа, AI загваруудын инференс (inference) болон хэрэглэгчийн мэдээллийн аюулгүй байдлыг хангах.
\end{itemize}

\section{Функционал шаардлага}
Системийн гүйцэтгэх үндсэн үйлдлүүдийг оролцогч тус бүрээр ангилав.

\textbf{Хэрэглэгчийн функциональ шаардлага:}
\begin{enumerate}
    \item Аудио файл (.wav, .mp3) системд байршуулж, CNN загвараар EQ/Compression анализ хийлгэх.
    \item MIDI файл дээр суурилсан MusicVAE ашиглан шинэ аяны хувилбар (variations) үүсгэх.
    \item RAG чатботаас техникийн болон онолын асуултад хариулт авах.
    \item Marketplace хэсгээс мэргэжлийн Sample Pack болон VST хэрэгслүүд хайх.
\end{enumerate}

\textbf{Админы функциональ шаардлага:}
\begin{enumerate}
    \item AI загваруудын сургалтын датасет болон Vector Database-ийг удирдах.
    \item Хэрэглэгчийн хандалт болон системийн ачааллын тайлан гаргах.
    \item Marketplace дээрх контентыг хянах, зөрчилтэй файлыг устгах.
\end{enumerate}

\section{Функционал бус шаардлага}
\begin{enumerate}
    \item \textbf{Хурд:} AI инференс хийх хугацаа (жишээ нь, аудио анализ) 10 секундээс хэтрэхгүй байх.
    \item \textbf{Сервер:} Нэгэн зэрэг 1000-аас дээш хэрэглэгч хандахад систем тогтвортой ажиллах (Scalability).
    \item \textbf{Найдвартай байдал:} RAG систем нь зөвхөн баталгаажсан эх сурвалжаас мэдээлэл хайж, Hallucination-ийг хамгийн бага түвшинд байлгах.
    \item \textbf{Стандарт:} Кодчилол нь Python PEP8 болон React стандартын дагуу хийгдэнэ.
\end{enumerate}

\section{Юзкейс (Use Case) диаграмын тодорхойлолт}

\begin{center}
    \begin{table}[!htbp]
        \caption{Юзкейс: Аудио микс анализ хийх}
        \begin{tabular}{|p{4cm}|p{11cm}|}
        \hline
            Нэр: & Аудио микс анализ хийх (AI Mixing Analysis) \\
        \hline
            ID: & UC-01 \\
        \hline
            Товч тайлбар: & Хэрэглэгч өөрийн аудио файлыг системд оруулж, AI-аас техникийн зөвлөмж авна. \\
        \hline
            Үндсэн оролцогч: & Хэрэглэгч \\
        \hline
            Туслах оролцогч: & AI Backend (CNN Model) \\
        \hline
            Өмнөх нөхцөл: & Хэрэглэгч системд нэвтэрсэн, аудио файл бэлэн байх. \\
        \hline
            Ажлын урсгал: & 
            \begin{enumerate}
                \item Хэрэглэгч аудио файл байршуулах хэсгийг сонгоно.
                \item Систем файлыг хүлээн авч, Спектрограмм (Spectrogram) болгон хувиргана.
                \item CNN загвар спектрограмм дээр анализ хийж, алдааг илрүүлнэ.
                \item Систем EQ болон Далайцын зөвлөмжийг хэрэглэгчид графиз хэлбэрээр харуулна.
            \end{enumerate} \\
        \hline
            Дараах нөхцөл: & Хэрэглэгч техникийн алдаагаа засах онолын зөвлөмжийг авсан байна. \\
        \hline
        \end{tabular}
    \end{table}
\end{center}



\section{Шинжилгээний класс болон дарааллын диаграмууд}

Системийн дотоод үйл ажиллагааг харуулахын тулд UML стандарт диаграмуудыг доор байршуулав.

\begin{figure}[htbp]
    \centering
    \includegraphics[scale=0.6]{Diagrams/Sequence_Analysis}
    \caption[Аудио анализын дарааллын диаграм]{Аудио анализын дарааллын диаграм}
    \label{fig:SeqAnalysis}
\end{figure}



\section{Бүлгийн дүгнэлт}
Хоёрдугаар бүлгийн хүрээнд системийн функционал болон функционал бус шаардлагуудыг тодорхойлж, хэрэглэгчдийн үйл ажиллагааг Юзкейс диаграммаар нарийвчлан гаргалаа. Энэхүү шинжилгээ нь дараагийн бүлэгт хийгдэх AI загваруудын интеграци болон өгөгдлийн сангийн дизайн хийхэд үндсэн удирдамж болох юм.